\section{Literature Review Plan}
\label{ch02:sec:plan}

The review was planned around the decision this research studies: what an agent should do with an IT incident whose cause is uncertain. It drew on two sources. The first was a core set of publications on AI for IT operations (AIOps), safe reinforcement learning and learning to defer, collected while the problem was being defined. The second was a systematic review, reported following the PRISMA 2020 guideline~\cite{page2021}, which searched Scopus, the ACM Digital Library and IEEE Xplore, with the Scopus and ACM searches limited to publications from 2015 to 2026, and included 34 studies that apply reinforcement learning or sequential decision-making to IT or security incident handling. The search strings are given in Appendix~\ref{App:scopus-search-strings} and the flow of records in Appendix~\ref{App:figures-other}.

The literature is discussed chronologically, in the order in which the strands developed: the foundations of safe reinforcement learning, then information gathering and deferral, then the prediction-based approaches of AIOps, then reinforcement learning for incident handling, and finally the recent studies that begin to combine these strands. For each approach the review states what it does, what evidence supports it and where it falls short, and then judges it against five requirements that an agent for IT incident triage at a university IT directorate must meet:

\begin{enumerate}[label=\textbf{R\arabic*.}]
\item \textbf{Sequential decisions:} the agent decides step by step, and each step can change what it knows.
\item \textbf{Costly investigation:} the agent can request a diagnostic check and weigh what the check costs against what it may reveal.
\item \textbf{Remedy or hand-over:} the agent can recommend a fix that changes the affected system, or hand the incident to a human expert.
\item \textbf{Operational rules:} every proposed action is checked against the organisation's rules, and a prohibited action is never carried out.
\item \textbf{Ambiguous evidence:} the agent stays reliable as the evidence becomes missing, unreliable or consistent with several causes.
\end{enumerate}

\section{Review of Approaches to Incident Triage}
\label{ch02:sec:review}

\subsection{Foundations of Safe Reinforcement Learning (2015--2018)}

These methods keep a learning agent from taking prohibited actions, and they are the source of the rule mechanism in this study (R4).

\begin{itemize}
\item \textbf{Constrained optimisation (2015--2017):} García and Fernández divided safe reinforcement learning into methods that change the optimisation criterion and methods that change exploration with external knowledge~\cite{garcia2015}. Constrained policy optimisation, from the first group, bounds the expected cumulative cost of a policy and guarantees near-constraint satisfaction at every training iteration~\cite{achiam2017}. Its strength is that safety is built into training rather than added afterwards. Its weakness is that the guarantee holds on average: an expected-cost bound does not enforce a rule that must hold at every step~\cite{zhao2023}, and a reward penalty discourages an action without preventing it~\cite{konighofer2025}. The rules of an IT directorate are not averages; a rule such as ``do not restart a production service without approval'' must hold on every incident, so this approach cannot meet R4 on its own.

\item \textbf{Shielding (2018 onwards):} Alshiekh et al.\ introduced a shield, a separate component that checks each action the agent proposes and blocks or replaces it when it would violate a specification, so that safety holds during both learning and deployment while the agent is restricted as little as possible~\cite{alshiekh2018}. For this study the main strength is the separation: the rules sit outside the learner, where IT staff can read and change them. Three weaknesses remain. The shield is computed from an abstraction of the environment's dynamics~\cite{alshiekh2018}, and shields that need a full model of the environment cannot be applied where no such model exists~\cite{yang2023}. A shield that silently replaces actions can hinder learning, because the agent does not know it was overridden~\cite{yang2023}. And shields have been demonstrated mainly in control tasks and games~\cite{alshiekh2018,konighofer2025}, not on the rules of IT operations. The rules of an IT directorate are conditions on the proposed action and the evidence gathered so far, so they can be checked without a model of how incidents evolve, but what happens to a blocked action has to be designed for this setting rather than borrowed.
\end{itemize}

Constrained optimisation shows that a penalty alone cannot meet R4, and shielding shows how R4 can be met. Neither addresses investigation, hand-over or ambiguous evidence.

\subsection{Information Gathering and Deferral (2018--2023)}

These methods let a model delay or hand over its decision, the two options that prediction-only triage lacks (R2 and R3).

\begin{itemize}
\item \textbf{Learning to defer (2018):} Madras et al.\ trained a model that can pass a decision to an external decision-maker and that accounts for the biases of that decision-maker; deferring made the combined system both more accurate and fairer, even when the human was inconsistent~\cite{madras2018}. Joshi et al.\ later extended deferral to sequential settings~\cite{joshi2023}, and selective prediction gives a model the related option to abstain when its confidence is low~\cite{geifman2019}. The strength of this line is that handing over is learned rather than set by a fixed threshold, and it is judged by how well the human and the model perform together. Its weakness for incident triage is that the model can only decide or hand over; it cannot ask for more evidence first. Madras et al.\ evaluated single classification decisions on recidivism and health datasets~\cite{madras2018}, where no further evidence can be collected. Applied to IT incidents, deferral alone would hand over incidents that one more check could have resolved.

\item \textbf{Information gathering (2019):} Janisch et al.\ trained an agent to acquire costly features one at a time before it classifies~\cite{janisch2019}, building on the principle that, under partial observability, an action may be chosen for the information it yields~\cite{kaelbling1998}. Its strength is that investigation is treated as a decision with a cost (R2), which is close to how IT staff work through diagnostic checks. Its weakness is the mirror image of deferral: the agent can look further or decide, but it cannot hand over, so when the evidence stays ambiguous it must still commit to an answer.
\end{itemize}

The two lines complement each other. Each supplies one of the options this study needs, neither supplies both, and neither checks its decisions against operational rules.

\subsection{Triage as Prediction in AIOps (2020--2025)}

These methods treat triage as the prediction of a label from the incident report.

\begin{itemize}
\item \textbf{Incident management at scale (2020):} Chen et al.\ conducted an empirical study of more than two years of incident management at Microsoft~\cite{chen2020}. Its strength is its evidence: drawn from real operations, it shows that distinct problems can produce similar symptoms, so that incidents are reassigned several times, which is the ambiguity this study simulates. The remedy it describes, however, is a recurrent model that predicts the responsible team from the incident text~\cite{chen2020}. Such a model makes one prediction per incident; when the text is ambiguous it still names a team, and the reassignments the study documents are the cost of naming the wrong one.

\item \textbf{AIOps surveys (2024--2025):} Remil et al.\ catalogued the methods for classifying, prioritising and assigning incidents, and criticised the reliance on classification metrics that ignore the cost of a wrong action~\cite{remil2024}. Zhang et al.\ observed that language-model assistants extended this work to diagnosis and mitigation suggestions, but that automation mainly stops at triage or recommendation~\cite{zhang2025}. The surveys are valuable as maps of the field and for the measurement problem they identify. The methods they survey, though, return an answer for every report and are scored on how often that answer is right. They cannot state that more evidence is needed or that a person should decide (R2, R3), and accuracy says nothing about whether a recommended action was permitted (R4).
\end{itemize}

This group supplies the problem and the evidence for it, but not a model of the decision. It corresponds to the single-prediction baseline against which this study's agent is compared.

\subsection{Reinforcement Learning for Incident Handling (2023--2026)}

These studies treat incident handling as a sequence of decisions (R1).

\begin{itemize}
\item \textbf{Root cause analysis (2023):} Wang et al.\ used hierarchical reinforcement learning to decide which questions to put to engineers when building a service dependency graph for Microsoft's microservice systems, reducing the number of queries needed from $O(N^2)$ to $O(1)$~\cite{wang2023hrlhf}. Ding et al.'s TraceDiag learned to prune the dependency graph before analysis and was deployed in production at Microsoft~\cite{ding2023}. Both treat investigation as a decision with a cost (R2), and both were built on real production systems, which few studies in this review achieve. Their goal, however, is to locate a root cause, not to decide what to do about it: neither recommends a remedy, checks one against rules, or decides whether a person should take over. Both also rely on the dependency graphs and traces of large microservice systems, whereas the Directorate diagnoses errors from API logs and a monitoring dashboard.

\item \textbf{Security incident response (2025--2026):} Dunsin et al.\ trained an agent to conduct a malware investigation step by step, aiming to shorten post-incident forensic investigations, and evaluated it in simulated scenarios~\cite{dunsin2025}. Hammar and Li planned incident responses while holding a belief over the model of the incident itself, relaxing the usual assumption that the system model is known, and bounded the errors this introduces~\cite{hammar2025}. Hammar et al.\ adapted response policies to changes in the system without repeating the offline optimisation~\cite{hammar2026belief}. The Hammar studies are the strongest in this group on R5: on the CAGE-2 benchmark, and for the 2026 study also on a testbed, they outperformed earlier methods in robustness to an inaccurate model~\cite{hammar2025,hammar2026belief}. Their weaknesses for this study are twofold. Their agents choose among responses or investigative steps, but none can hand the incident to a person (R3), and none checks its actions against an organisation's rules (R4). Their environments also model cyberattacks, in which an adversary drives the incident, whereas the IT service incidents in this study arise from faults and misconfigurations.
\end{itemize}

This group shows that reinforcement learning can handle investigation and response as sequential decisions. In every study, however, the decision ends with the agent: none passes an incident to a person or respects operational rules.

\subsection{Recent Combinations (2025--2026)}

The most recent studies begin to combine the strands above, mostly in security operations.

\begin{itemize}
\item \textbf{Gathering evidence and escalating in alert triage:} Okafor trained offline agents that dismiss, investigate or contain security alerts, with escalation to analysts fixed by design~\cite{okafor2026}. Its strength is methodological: it avoids unsafe online exploration, removes temporal leakage from the evaluation, and shows that F1 and ROC-AUC can hide degenerate policies~\cite{okafor2026}. Its weakness is its data: public intrusion datasets hold flow records rather than decisions, so its training episodes are reconstructed from flows~\cite{okafor2026}, and the escalation it reports was not learned. Sahin and Mert trained an agent that queries logs and threat intelligence before it escalates or closes an alert, with a separate conformal layer that defers uncertain alerts to an analyst while controlling the rate of missed threats~\cite{sahin2026}. This gives hand-over a statistical guarantee rather than a tuned threshold, the strongest treatment of R3 in the review, but the deferral sits outside the learned policy: the agent itself is not taught when handing over is worth its cost. Kovtun et al.\ learned a policy that chooses among shallow triage, deep analysis and deferral under analyst-capacity constraints, and theirs is the one study that reports decision quality, cost, escalation and compliance together~\cite{kovtun2026}. Its evaluation, like Okafor's, rests on public intrusion datasets, CIC-IDS2017 and UNSW-NB15~\cite{kovtun2026}. Across the three, the actions close, contain or escalate an alert; none recommends a repair to a failing service, which is the action an IT directorate most needs to check against its rules.

\item \textbf{Rule mechanisms in incident response:} Ravi et al.'s SAFE-HGX added a safety layer to a hierarchical multi-agent learner for intrusion response in industrial control systems, replacing blocked actions, and in simulated attack scenarios it learned faster and contained attacks more accurately than DQN and PPO baselines~\cite{ravi2025}. The cost was high: about 9 to 11 times as many false-positive actions as the PPO baseline, and 74 to 77\% higher operating cost~\cite{ravi2025}. An IT directorate is unlikely to accept that trade-off, since an unnecessary action on a live service can itself disrupt operations~\cite{okafor2026}. Singh's deterministic guard for cloud remediation hands an action to a human when it fails~\cite{singh2026}. In its ablation, removing the guard left the mean time to resolve almost unchanged (144~s against 142~s) but cut governance compliance from 99.9\% to 72.3\%~\cite{singh2026}, the clearest evidence in the review that a hard rule check need not slow resolution. The learning in Singh's framework, however, is confined to a reinforcement learning meta-orchestrator that weights the recommendations of other agents~\cite{singh2026}, so the study does not show how a learning agent behaves when its own proposals are blocked. Alauthman et al.\ proposed a shielded DQN agent with operator override for incident response in industrial control systems, but reported an evaluation plan and indicative findings rather than measured results~\cite{alauthman2026}; it shows that the design is plausible, not that it works. Of the 34 studies in the systematic review, only three say what happens to a blocked decision, and none lets the agent propose again.
\end{itemize}

These two groups have reached the design questions of this study from opposite sides. The alert-triage studies hand cases over but do not check remedies against rules; the rule-mechanism studies check actions against rules but do not learn when to hand over. Their results also disagree on what a rule mechanism costs: Ravi et al.\ found many more false actions, while Singh found no loss of speed.

\subsection{Evaluation and Robustness}

These studies show how such agents have been evaluated, and where the evaluations fall short.

\begin{itemize}
\item \textbf{Measures:} In the systematic review, 31 of 34 studies measure decision quality, but only five measure escalation behaviour and four constraint compliance; one study, Kovtun et al., reports all four kinds of measure, for SOC alert escalation~\cite{kovtun2026}. This matters because each measure on its own can be satisfied by a poor policy: an agent that escalates every incident never breaks a rule and never recommends a wrong fix, but it is of no use to the people it escalates to. A single metric can also hide a degenerate policy: Okafor found that F1 and ROC-AUC concealed agents that had collapsed into trivial detectors, which the Matthews correlation coefficient exposed~\cite{okafor2026}.
\item \textbf{Robustness:} Fifteen of the reviewed studies test their agents under changed conditions, and the methods that keep a belief over the hidden state degrade least~\cite{hammar2025,hammar2026belief}. Sahin and Mert show that deferral to analysts rises as the evidence degrades~\cite{sahin2026}. These tests vary the conditions as a whole; no reviewed study reports how the mix of investigating, acting and escalating changes as ambiguity increases, which is the behaviour that shows whether an agent recognises when it does not know enough.
\end{itemize}

Table~\ref{tab:comparison} compares the studies closest to this research against the capabilities it requires.

\begin{table}[H]
\centering
\caption{Comparison of the studies closest to this research. \cellyes~=~addressed; \cellpart~=~partly addressed; \cellno~=~not addressed.}
\label{tab:comparison}
\small
\setlength{\tabcolsep}{4pt}
\begin{tabular}{|p{3.4cm}|p{4.2cm}|c|c|c|c|c|}
\hline
\textbf{Study} & \textbf{Setting} & \textbf{Seq.} & \textbf{Info.} & \textbf{Rule} & \textbf{Esc.} & \textbf{Amb.} \\
\hline
Chen et al.~\cite{chen2020} & IT incident triage & \cellno & \cellno & \cellno & \cellno & \cellno \\ \hline
Wang et al.~\cite{wang2023hrlhf} & Root cause analysis & \cellyes & \cellyes & \cellno & \cellpart & \cellno \\ \hline
Dunsin et al.~\cite{dunsin2025} & Malware investigation & \cellyes & \cellyes & \cellno & \cellno & \cellno \\ \hline
Hammar and Li~\cite{hammar2025} & Security incident response & \cellyes & \cellno & \cellno & \cellno & \cellyes \\ \hline
Alauthman et al.~\cite{alauthman2026} & Industrial control incident response & \cellyes & \cellno & \cellyes & \cellpart & \cellno \\ \hline
Ravi et al.~\cite{ravi2025} & Industrial control incident response & \cellyes & \cellno & \cellyes & \cellno & \cellno \\ \hline
Okafor~\cite{okafor2026} & Security alert triage & \cellyes & \cellyes & \cellpart & \cellpart & \cellno \\ \hline
Sahin and Mert~\cite{sahin2026} & Security alert triage & \cellyes & \cellyes & \cellpart & \cellyes & \cellyes \\ \hline
Kovtun et al.~\cite{kovtun2026} & Security alert escalation & \cellyes & \cellpart & \cellpart & \cellyes & \cellno \\ \hline
Singh~\cite{singh2026} & Cloud operations remediation & \cellpart & \cellno & \cellyes & \cellpart & \cellno \\ \hline
\textbf{This study} & IT incident triage & \cellyes & \cellyes & \cellyes & \cellyes & \cellyes \\ \hline
\end{tabular}

\vspace{4pt}
\raggedright\footnotesize Seq.\ = sequential policy; Info.\ = information gathering; Rule = rule check; Esc.\ = human escalation; Amb.\ = varied ambiguity.
\end{table}

\section{Research Gaps}
\label{ch02:sec:gaps}

The approaches reviewed above each solve part of the triage decision, but no reviewed study combines them for IT service incidents. The most significant gap is in the decision itself: information-gathering studies let an agent choose between looking further and deciding, deferral studies let it choose between deciding and handing over, and the security alert-triage studies that join the two do not recommend a fix that acts on the system. No reviewed study trains one policy to choose among investigating, recommending a remedy and escalating for IT service incidents. The second gap concerns the rules: hard rule mechanisms have been measured in only two studies, their results point in opposite directions, and no reviewed study gives a blocked agent a second proposal or reports how often blocked proposals occur during training. The third gap concerns robustness: decision accuracy and the rate of deferral have been tested under degraded evidence, but not the mix of investigating, recommending and escalating as ambiguity rises. Underneath all three lies a gap in representation, since the available environments model cyberattacks and microservice faults, and none lets the ambiguity of IT service incidents be varied.

These gaps are stated for the literature reviewed. The systematic review searched three databases, limited its IEEE Xplore search to open-access records, did not carry out citation searching, and could not retrieve seven reports. The claim is that the reviewed literature leaves these questions open, not that they have never been examined.
